\section{Week 9 - 9/11}
The goal of this week is to get some things written down. I've made a list of things that need to be explained, and I'd like to figure out what the best structure is for these things
\begin{itemize}
    \item Green's functions, and their relations to observables.
    \begin{itemize}
        \item Maybe look into high energy connections?
        \item Some kind of derivation of Keldysh.
        \item Why does the Green's function relate to DOS?
    \end{itemize}
    \item The basic model, its derivation, and the result for the current.
    \item Why we expect generality.
    \begin{itemize}
        \item Maybe show the other derivations.
        \item Dive into \underline{coherent} versus \underline{incoherent}.
    \end{itemize}
    \item The noise derivation.
    \begin{itemize}
        \item Immediately do the resummation.
        \item Maybe do the derivation at zero energy.
        \item Kramers-Kronig and the large bandwidth approximation.
    \end{itemize}
    \item Energy dependence of the noise.
    \begin{itemize}
        \item Why only the zero frequency value is important.
        \item Refer to Ingmar's paper, and the experimental unreachability of the high frequency regime.
    \end{itemize}
    \item Explain why this is relevant to begin with.
    \item Splitting of the noise into thermal and quantum noise, and why we chose what we chose.
\end{itemize}

I want to subdivide it into chapters. The distribution of chapters should be something like

\begin{itemize}
    \item Relevance
    \item Green's functions, linear response theory, Wiener-Khinchin, non-equilibrium.
    \item Normal metal current and generality.
    \item Normal metal noise, separated into energy dependence and resummation
\end{itemize}